% page 338 of the facsimile (image p-337 of the scan)
% block 147.3 x 224.7 mm ; left margin 8.0 mm
\pgc{-12.700mm}{0.000mm}{
\l{\cel{3}330.}
\l{}
\l{\sou{lanugo}. I. Mikra plumi dina e dolca qui kreskas (en la}
\l{\cel{3}komenco di lua vivo) che la ucelo e formacas strato sub}
\l{\cel{3}la plumi granda. - II. Mikra pili quin havas, lor sua}
\l{\cel{3}nasko, la maxim multa quadripedi, e qui duras garnisar}
\l{\cel{3}la korpo sub la pili plu forta qui formacas la pilaro.3}
\l{\cel{3}Sorto di kotono lejera qua tegas la stipo, la squami}
\l{\cel{3}di la burjoneti, la folio, la frukto, di ula vejetanti.}
\l{\cel{3}- IV. Barbo naskanta, pilo tre tenua. - EFI.}
\l{}
\l{\sou{lapar}.\cel{2}(\sou{trans}.) (\sou{aludante precipue la hundi}). Drinkar per}
\l{\cel{3}quaze pumpar la liquido per sua lango. - DEF.}
\l{}
\l{\sou{lapido}. Stono tre kustoza, kom ex.: rubino, agato, onixo,}
\l{\cel{3}kornalino, esmeraldo, safiro, e c. - DEFS.}
\l{}
\l{\sou{lapislazulo}. (\sou{minera}l.) Stono opaka, blua, kun veini blan-}
\l{\cel{3}ka, e kun punti ek piriti feroza qui cintilifas quale}
\l{\cel{3}oro. DEFIRS.}
\l{}
\l{\sou{lardo}. Graso qua garnisas la pelo di la porko e di kelka}
\l{\cel{3}altra animali, precipue ye la ventro. - FIS.}
\l{}
\l{\sou{larico}. (\sou{bot.}) Arboro konifera, de la trunko di qua de-}
\l{\cel{3}kulas rezino e di qua la folii sekrecas substanco vis-}
\l{\cel{3}kozo. - L.\cel{1}\sou{larix}. - DEFISL.}
\l{}
\l{\sou{laringo}.\cel{2}(\sou{anat.}) Organo esencala di voco, formacita ek}
\l{\cel{3}plura peci intermoviva, ye la parto supra di la tra-}
\l{\cel{3}keo.-DEFIS.}
\l{}
\l{\sou{laringito}. (\sou{patol}.)Inflamuro di la laringo. - DEFIS.}
\l{}
\l{\sou{laringoskopo}. (\sou{tek}.)Aparato per qua onu povas observar la}
\l{\cel{3}parto interna di la laringo. - DEFS.}
\l{}
\l{\sou{laringotomio}. (\sou{kirurg}.) Operaco per qua onu facas aperturo}
\l{\cel{3}en la laringo por preventar asfixio. - DEFIS.}
\l{}
\l{\sou{larja}. Tre extensita en la sinso opozata a longeso. - FI.}
\l{}
\l{\sou{larvo}. I. (\sou{epoki antiqua, en Roma}).Fantomo ledega, repu-}
\l{\cel{3}gnanta. - II. (\sou{historio naturala)}. Insekto vermoforma}
\l{\cel{3}qua reprezentas la stando unesma di la metamorfos-}
\l{\cel{3}insekti. - DEFIS.}
\l{}
\l{\sou{lasar}.\cel{2}(\sou{trans}.) I. Igar ulo laxa per cesar retenar lu.}
\l{\cel{3}- II. Igar (ulu od ulo) ube lu esas - ne prenar lu kun}
\l{\cel{3}su. - III. Igar ulo,gardata,- ne deprenar lu de ulu.}
\l{\cel{3}- IV. Igar tale ke ulu posedeskas ulo quon onu cesas}
\l{\cel{3}gardar. - DFI.}
}
